% Job posting: https://ca.linkedin.com/jobs/view/r-d-biomedical-engineer-on-site-at-myant-4463171428
\documentclass[11pt]{article}
\usepackage[left=1in,top=1in,right=1in,bottom=1in]{geometry}
\usepackage[hidelinks]{hyperref}
\usepackage{setspace}
\usepackage[T1]{fontenc}
\usepackage{newtxtext,newtxmath}
\ifdefined\pdfgentounicode
  \input{glyphtounicode}
  \pdfgentounicode=1
\fi
\setstretch{1.1}
\pagenumbering{gobble}
\begin{document}
\pagestyle{empty}

\noindent
Nishal Stanislaus Silva, PhD \\
Toronto, ON \\
nishal.silva@hotmail.com \, | \, +1 613 200 2030 \\
nishal.xyz \, | \, linkedin.com/in/nishal-silva \, | \, github.com/ns2max

\vspace{1em}
\noindent Dear Hiring Manager,

\vspace{0.5em}
\noindent I am writing to apply for the R\&D Biomedical Engineer role at Myant. My background sits on an overlap I rarely see asked for on one job posting, and which maps directly onto textile-integrated biosensing: 5.5 years inside apparel/textile manufacturing at MAS Holdings, South Asia's largest apparel manufacturer, where I built end-to-end computer-vision and IoT systems on live production lines; hands-on wearable IMU sensor integration and power-budget engineering on a self-powered, battery-constrained embedded instrument at McGill University's IDMIL; and a real clinical health-monitoring deployment -- a computer-vision device I built at MAS that digitized patient vitals and was deployed across hospital wards during COVID-19, recognized by the Government Medical Officers' Association (GMOA), Sri Lanka. Textile manufacturing floor experience, wearable sensor hardware, and a shipped clinical monitoring device rarely sit in one CV, and together they speak directly to the engineering problem behind a textile-integrated biosensor.

\vspace{0.5em}
\noindent My PhD and postdoctoral work at the University of Trento extended this into biosignal hardware and embedded systems discipline: I built a real-time multimodal synchronization pipeline fusing audio, BCI/EEG biosignals (g.tec hardware), and mixed-reality head/hand tracking across four performers simultaneously, and I shipped a real-time pattern detector running in 14\,ms on a Raspberry Pi 4 at F1 0.76, 74x faster than the DTW baseline it replaced. At McGill, I integrated a BNO055 IMU over I2C into a self-powered instrument and profiled its compute and power draw to validate real-time feasibility on battery-constrained hardware -- the same power-budget discipline a wearable biosensor demands. My core engineering languages are C++17 and Python, applied throughout to production and research systems alike.

\vspace{0.5em}
\noindent I have not manufactured e-textiles myself -- conductive yarn integration or textile-electrode fabrication is not on my CV -- and I say that plainly rather than gloss over it. What I bring instead is the rarer combination: years of manufacturing-floor fluency in the textile industry itself, wearable inertial-sensor integration and power-budget engineering, and a computer-vision health-monitoring device that actually reached hospital wards. Few candidates for this role will have stood on a textile production floor, wired up a wearable IMU, and shipped a clinical monitoring device -- I have done all three, and I am confident that combination transfers directly to Myant's R\&D problems.

\vspace{0.5em}
\noindent I am a Canadian Permanent Resident, eligible to work in Canada without sponsorship, based in Toronto -- a short commute to Myant's Mississauga site -- and available to start immediately. I would welcome the opportunity to discuss how my background can contribute to your R\&D team.

\vspace{1em}
\noindent Sincerely, \\
Nishal Stanislaus Silva

\end{document}
